% ============================================================
% MANUSCRITO — IEEE JTEHM (Journal of Translational Engineering in
% Health and Medicine)
% Revisado 03-ago-2026 CONTRA LA PLANTILLA OFICIAL REAL que el usuario
% descargó de IEEE (docs/manuscrito/JTEHM_LaTex_Template/JERM Demo.tex +
% IEEEJERM.cls) — reemplaza la versión anterior de este archivo, que
% usaba IEEEtran genérico sin poder verificar la plantilla real.
%
% ============================================================
% ESTE ES EL ARCHIVO FINAL — el espejo de lo que está en Overleaf.
% Reubicado aquí el 05-ago-2026 (antes: docs/manuscrito/
% plantilla_overleaf_JTEHM.tex) para que viva junto al .cls y al .bib
% que necesita para compilar. Ver README_ESTRUCTURA.md en esta carpeta.
%
% Los BORRADORES en prosa siguen en ../metodos_introduccion_borrador.md
% — se redacta allá primero, y solo lo que ya está resuelto pasa a este
% archivo. Si los dos difieren, ESTE archivo manda (es lo que se envía).
% ============================================================
%
% CÓMO USAR ESTE ARCHIVO EN TU OVERLEAF:
% 1. En tu proyecto de Overleaf (el que ya creaste con el ZIP de IEEE),
%    reemplaza el contenido de "JERM Demo.tex" por el contenido de este
%    archivo completo (desde \documentclass hasta \end{document}).
% 2. NO borres del proyecto: IEEEJERM.cls, references.bib — sin el
%    primero no compila nada (\documentclass falla), sin el segundo
%    fallan las citas. Subir el references.bib DE ESTA CARPETA, que ya
%    tiene las referencias reales del manuscrito (el original del ZIP
%    solo traía dos entradas de ejemplo).
% 2-bis. Claves de cita actualizadas el 05-ago-2026: las
%    \cite{PENDIENTE_*} se reemplazaron por las claves reales.
% 2-ter. ACTUALIZADO 11-ago-2026: las tres entradas que estaban
%    incompletas ya están verificadas y completas, y por eso CAMBIARON
%    DE CLAVE (ahora Autor+Año, como el resto):
%      CompactGaitSimulator2024    -> Sudeesh2024CompactGaitSimulator
%      ConceptualDesignSimulator2015 -> Marinelli2015ConceptualDesign
%      IMU2023TranstibialValidity  -> Rattanakoch2023TranstibialIMU
%    Si pegaste una versión anterior en Overleaf, reemplaza también el
%    references.bib o las citas quedarán en [?].
%    Las 5 referencias de simuladores/instrumentación están [OK]; las 4
%    de métodos estadísticos siguen [REVISAR]. Ver ../referencias_verificadas.md
% 3. Cambios reales encontrados al comparar con la plantilla oficial
%    (correcciones respecto a la versión anterior de este archivo):
%    - Clase real: IEEEJERM (NO "IEEEtran" genérico), con opciones
%      [journal,twocolumn,letterpaper].
%    - El abstract tiene un campo QUE NO SABÍAMOS: "Clinical impact:"
%      además de Objective/Methods and procedures/Results/Conclusion.
%    - El "Clinical and Translational Impact Statement" (≤30 palabras,
%      obligatorio) además debe indicar una categoría del espectro
%      clínico del NIH: Early/Pre-Clinical Research, Clinical Research,
%      Clinical Implementation, o Public Health — esto NO estaba en el
%      resumen de la página pública de instrucciones, solo aparece en
%      la plantilla real.
%    - SÍ HAY sección "Discussion" separada de "Conclusion" — la
%      página pública de instrucciones decía que no la había para
%      "Papers"; la plantilla real la trae explícita, así que manda
%      esta. Se corrigió la estructura respecto a la versión anterior.
%    - Bibliografía: \bibliographystyle{ieeetr} (no "IEEEtran" a secas)
%      + \bibliography{references}.
%    - Subsecciones de Metodología en el demo usan \subsection*{...}
%      (sin numerar) — se mantiene esa convención aquí.
% 4. Sigue pendiente sin verificar: el límite exacto de páginas (8,
%    verificado en la página pública, no contradicho por la plantilla)
%    y el orden alfabético de keywords (tampoco contradicho).
% ============================================================

\documentclass[journal,twocolumn,letterpaper]{IEEEJERM}

\usepackage{times,amsmath,epsfig}
\usepackage{fancyhdr}
\usepackage{amsfonts}
\usepackage{amssymb}
\usepackage{array}
\usepackage{graphicx}
\usepackage{url}
\usepackage{subfigure}
\usepackage{bm}
\usepackage{breqn}
\usepackage{xcolor}
\usepackage{soul}
\usepackage[breaklinks=true]{hyperref}
\usepackage{breakcites}

\begin{document}

\vspace{1 cm}

\title{Multi-Subject Functional Validation of a 3-DOF Gait Simulator for
Transtibial Prosthesis Testing Using Inertial Motion Capture}

% [PENDIENTE: nombres, afiliaciones y orden de autoria. Borrador de
% roles CRediT ya armado en docs/manuscrito/creditos_autoria_borrador.md
% (13-ago-2026) -- confirmar primero si JTEHM pide CRediT en el .tex o
% en el portal de envio antes de agregar una seccion nueva aqui.]
\author{Autor~Uno$^1$,
        Autor~Dos$^1$,
        Autor~Tres$^2$,
        Autor~Cuatro$^{3}$
}

\markboth{IEEE Journal of Translational Engineering in Health and Medicine}
{Autor Uno \MakeLowercase{\textit{et al.}}: Multi-Subject Functional Validation of a 3-DOF Gait Simulator}

\twocolumn[
\begin{@twocolumnfalse}

\maketitle

\begin{abstract}

Objective: \textnormal{[PENDIENTE].}
Methods and procedures: \textnormal{[PENDIENTE].}
Results: \textnormal{[PENDIENTE -- depende de datos de sujetos nuevos y
del piloto del IMU en la plataforma; ver CLAUDE.md. Depende tambien de
que la integracion Raspberry Pi-ESP32 se resuelva -- sin eso no hay
salida del simulador que medir, ni cinematica ni cinetica.]}
Conclusion: \textnormal{[PENDIENTE].}

Clinical impact: \textnormal{[PENDIENTE -- una o dos frases sobre la
relevancia clinica practica, distinto del Impact Statement de abajo. NO
enmarcar como "medimos angulos con precision" -- JTEHM prioriza TRL 5-9
(tecnologia validada/demostrada en entorno relevante) y rechaza
explicitamente "lab validation without clinical context" (TRL 1-4).
Enmarcar como: esta herramienta reduce el tiempo/costo/riesgo de pruebas
continuas con sujetos humanos durante el desarrollo de protesis,
acelerando la iteracion de diseno antes de que un paciente real la use
-- calza con el alcance oficial "Technological relevance to healthcare
cost reduction" y "Technology affecting healthcare management". Ver
docs/manuscrito/guia_autor_JTEHM.md seccion 1-ter para el detalle
completo de este hallazgo.]}

% [RESUELTO 13-ago-2026: Impact Statement aprobado por el usuario en
% DISCUSION_Q2.md T-1, con una correccion -- se quito la mencion a
% "knee" (el simulador es transtibial, nunca cruza la rodilla). 25
% palabras, bajo el limite de 30. Categoria NIH: Early/Pre-Clinical
% Research, ya confirmada como una de las 4 categorias oficiales de la
% plantilla real de JTEHM.]
Clinical and Translational Impact Statement:
\textnormal{Early/Pre-Clinical Research. A validated bench simulator
lets transtibial prosthetic components be iterated and evaluated
without the repeated, fatiguing human-subject sessions that gait-lab
testing otherwise requires.}
\end{abstract}

\begin{IEEEkeywords}
gait simulator, inertial motion capture, multi-subject validation, prosthetics testing, transtibial prosthesis
\end{IEEEkeywords}
% [PENDIENTE: confirmar limite maximo de keywords y si deben ir en
% orden alfabetico -- la pagina publica de instrucciones decia hasta 5
% en orden alfabetico; la plantilla real no lo contradice ni lo confirma
% explicitamente, ya estan en ese orden arriba por las dudas]

\end{@twocolumnfalse}]

% Pie de pagina con afiliaciones -- ajustar segun autores reales
{
  \renewcommand{\thefootnote}{}%
  % [RESUELTO 13-ago-2026: correo de correspondencia confirmado en DISCUSION_Q2.md P-8 -- Luis Marcos Plasencia Janampa, autor de correspondencia]
  \footnotetext[1]{Departamento de Ingenier\'ia, Pontificia Universidad Cat\'olica del Per\'u (PUCP), Lima, Per\'u. Corresponding author: L. M. Plasencia Janampa (e-mail: luis.plasencia@pucp.edu.pe).}
  \footnotetext[2]{[PENDIENTE: afiliaci\'on]}
  \footnotetext[3]{[PENDIENTE: afiliaci\'on]}
}

\IEEEpeerreviewmaketitle

% ============================================================
\section{Introduction}
% ============================================================
% [RESUELTO 11-ago-2026: las 4 citas de este primer parrafo ya estan
% verificadas contra la fuente y completas en references.bib.
% Queda un pendiente menor: DeRaeve2014AlignmentTool es un resumen de
% congreso de una pagina (cita debil para IEEE) y sigue sin reemplazo
% identificado -- ver nota en references.bib.]
% [PENDIENTE: RECORDAR limite de 8 paginas totales (verificado en la
% pagina publica de instrucciones) -- esta Introduccion es la primera
% candidata a recortar si el manuscrito se pasa de longitud.]

Robotic and mechanical gait simulators are an established approach for
prosthetic device development and testing, offering repeatable dynamic
loading conditions that are difficult to achieve through human-subject
trials alone, while reducing the logistical and ethical burden of
continuous human involvement and improving repeatability across
prototype comparisons \cite{Etoundi2022RoboticTestRig}. Cost-effective
gait simulators in particular have been proposed specifically to reduce
dependency on user trials during prosthesis development
\cite{Sudeesh2024CompactGaitSimulator}. Existing simulators for
lower-limb prostheses range from simplified sagittal-plane rigs
reproducing hip-to-foot vertical displacement and knee flexion-extension
\cite{Marinelli2015ConceptualDesign} to robotic platforms used for
prosthetic alignment and joint performance assessment
\cite{DeRaeve2014AlignmentTool,Etoundi2022RoboticTestRig}.
% [RESUELTO 13-ago-2026, ver DISCUSION_Q2.md P-11: busqueda exhaustiva no
% encontro reemplazo de articulo completo para DeRaeve2014AlignmentTool --
% el grupo De Raeve/Muraru/Peeraer solo publico el resumen de congreso de
% una pagina para este sistema exacto, confirmado por multiples busquedas
% independientes. Se decidio NO inventar ni forzar una cita tangencial;
% en su lugar se cito junto a Etoundi2022RoboticTestRig (ya citado arriba
% en el mismo parrafo, articulo completo revisado por pares en Frontiers
% in Robotics and AI) para reforzar la mitad de la afirmacion que cubre
% "joint performance assessment". DeRaeve sigue siendo la unica fuente
% real para la mitad de "alignment" -- se mantiene porque describe
% correctamente lo que es, no se elimina la clausula.]

However, the validity of any such simulator rests on a largely untested
assumption: that a trajectory programmed from a single reference
subject's gait is representative of the broader population it is
intended to serve. This question is rarely addressed with quantified,
multi-subject evidence in the prosthesis gait simulator literature. The
present study addresses this gap directly: we assess whether a
3-degree-of-freedom (horizontal, vertical, sagittal) motor-driven gait
simulator reproduces gait kinematics captured from multiple subjects
independent of its calibration trajectory, using an inertial motion
capture system previously validated against an optoelectronic reference
for sagittal-plane lower-limb kinematics
\cite{Piche2022iSenValidity,Rattanakoch2023TranstibialIMU}, and we quantify ---
rather than narrate --- the sources of vertical force overestimation
inherent to a motor-driven platform. By replicating gait kinematics
captured from human subjects who did not participate in its
programming, the simulator functions as a pre-clinical evaluation
bench at TRL 5--6, positioned to accelerate prosthetic component
iteration ahead of, and in support of, in-clinic assessment.
% [RESUELTO 13-ago-2026: frase de posicionamiento TRL aprobada por el
% usuario en DISCUSION_Q2.md T-1, sin cambios sobre el borrador.]

% ============================================================
\section{Methodology}
% ============================================================

\subsection*{System Description}
% Condensado, autocontenido -- no remite al paper de conferencia (no citado en este articulo)

The gait simulator reproduces a pre-recorded gait trajectory across
three independent degrees of freedom --- horizontal displacement,
vertical displacement, and sagittal rotation --- each driven by its own
motor and transmission (chain-and-sprocket or direct drive, depending on
the axis). This sagittal-plane architecture is consistent with an
independent analysis of the minimum degrees of freedom required for
prosthetic testing, which concluded that flexion-extension together with
vertical and horizontal translation are sufficient
\cite{Sudeesh2024CompactGaitSimulator}. Trajectories are stored as CSV files, uploaded via a
Raspberry Pi, and executed open-loop by an ESP32 microcontroller; the
present study does not modify this operating principle --- no
closed-loop control or algorithmic trajectory generation was introduced
in this cycle (see Discussion, Future Work). % [PENDIENTE: calibracion del actuador vertical, bloqueada por integracion Raspberry Pi-ESP32]

\subsection*{Instrumentation}
% Un solo instrumento en todo el protocolo (revisado 03-ago-2026)

Kinematic and kinetic data were collected using two instruments: an
inertial motion capture system (STT-IWS/iSen) and a force platform
(AMTI). The inertial system was used throughout the protocol to (i)
capture new subjects' natural gait, (ii) re-capture the original
reference subject to regenerate the simulator's default trajectory, and
(iii) measure the simulator's own platform output (sensor mounted
directly on the platform). Segment/platform inclination angles are
derived from the sensor's onboard orientation estimate relative to
gravity, positive above horizontal and negative below, the same sign
convention used throughout this project. No in-house cross-instrument
validation was performed in this cycle; instrument validity was taken
from published evidence rather than re-derived. The specific system used
here has been validated against an optoelectronic reference (OptiTrack)
for sagittal-plane lower-limb kinematics, reporting waveform RMSD of
3.3$^\circ$ for the knee and 7.3$^\circ$ for the hip, Lin's concordance
correlation above 0.75 for all joints except the ankle, and excellent
test-retest repeatability (ICC 0.93 knee, 0.89 hip)
\cite{Piche2022iSenValidity}. Sensor drift over a three-minute
acquisition was negligible in that study (maximum relative error
1.1$^\circ$; RMSD below 2$^\circ$ between the first and last 30\,s).
Validity of inertial motion capture in the population of interest is
supported separately, in transtibial prosthesis users, by
\cite{Rattanakoch2023TranstibialIMU}.
% NOTA 11-ago-2026 -- POR QUE ESTAN LAS DOS CITAS Y NO UNA:
% Piche2022 valida EL INSTRUMENTO de este proyecto (iSen STT-IWS), pero
% en 22 sanos jovenes. Rattanakoch2023 cubre la POBLACION transtibial,
% pero con OTRO instrumento (Noraxon MyoMotion). Ninguna cubre las dos
% cosas; por eso se presentan como complementarias y de forma explicita.
% NO fusionarlas ni escribir que una valido el instrumento en poblacion
% transtibial -- seria falso y trivialmente verificable.
% No se cita el LCC de tobillo (0.72/0.60/0.56, cae con la velocidad)
% porque este estudio no mide angulo de tobillo. Es defendible, pero hay
% que saber sostenerlo si un revisor pregunta.
The
force platform recorded vertical ground reaction force (Fz) exerted by
the simulator against a fixed load cell, sampled at 1000\,Hz.
% [PENDIENTE: confirmar orientacion cruda vs. angulo articular calibrado
% -- resultado del piloto en curso. Ver docs/literatura/postprocesado_datos_crudos_IMU.md.]

\subsection*{Vertical Force Overestimation: A Three-Stage, Decoupled Explanation}

Because the moving assembly is not a single rigid body --- a stationary
electronic control board is mounted on one side of the frame, while
three independently actuated axes (horizontal, vertical, sagittal) each
carry their own kinematic chain on the other --- vertical force
overestimation is decomposed into three independent, sequentially
corrected sources rather than a single inertial term: (1) a fixed
vertical datum offset, calibrated once against an independent static
load reference and held constant across all trials; (2) command-to-encoder
tracking fidelity per axis, isolating mechanical/control error from the
remaining sources; (3) an axis-wise inertial correction, using the mass
that actually accelerates with each motor according to the real
mounting order of the kinematic chain, benchmarked against published
prosthetic gait GRF data.
% [PENDIENTE: fuente de m_eje sin resolver -- CAD vs. renunciar al termino y declararlo limitacion]

\subsection*{Experimental Protocol}
% [PENDIENTE: depende de si llego la aprobacion de etica y del numero final de sujetos]
Please insert text here.

\subsection*{Statistical Analysis}

Kinematic and kinetic curves were compared using the normalized
root-mean-square error (RMSEnorm), computed from the point-wise
difference between simulator and reference signals normalized by the
reference's point-wise standard deviation. Pearson's correlation
coefficient ($r$) was reported alongside RMSEnorm but interpreted with
caution during monotonic phases of the gait cycle, where $r$ is known to
be insensitive to offset and scaling errors. The percentage of points
falling within $\pm$1 SD of the reference curve was reported as an
intuitive tolerance envelope. Trial-to-trial repeatability was
quantified using the intraclass correlation coefficient ICC(3,1)
(two-way mixed-effects model, absolute agreement, single measurement)
\cite{KooLi2016}. To localize differences across the gait cycle rather
than relying solely on scalar summaries, non-parametric Statistical
Parametric Mapping (SPM1D) was applied via permutation (sign-flip for
paired designs, label-shuffling for independent-group designs), with the
family-wise error rate controlled through the maximum-statistic
threshold \cite{NicholsHolmes2002}. The non-parametric variant was
preferred over classical random-field-theory SPM given the limited
trial counts expected in this study (5--10 trials per subject, up to
15--20 subjects), consistent with recommendations for small-sample
biomechanical datasets \cite{Pataky2015}. This trueness/precision
framework follows ISO 5725 terminology \cite{ISO5725}, with RMSEnorm and
CMC representing trueness (closeness to a reference value) and
ICC(3,1)/between-subject coefficient of variation representing precision
(repeatability).
% [PENDIENTE: revisar longitud contra el limite de 8 paginas una vez
% que el resto del manuscrito este redactado -- recortar si hace falta.
% Confirmar claves de cita exactas (KooLi2016, NicholsHolmes2002,
% Pataky2015, ISO5725) contra references.bib cuando se arme.]

% ============================================================
\section{Results}
% ============================================================
% [PENDIENTE: depende de datos de sujetos nuevos y de que se resuelva
% la integracion Raspberry Pi-ESP32 -- sin eso no hay salida del
% simulador que medir (ni cinematica ni cinetica), ver CLAUDE.md]

\subsection*{Tracking Fidelity per Subject}
% [RESUELTO 13-ago-2026, esqueleto: tabla y figura listas, sin valores
% -- columnas tomadas directamente de resultado.por_sujeto en
% CODIGOS/MULTISUJETO/Procesar_Multisujeto_Core.m (ver
% GUIA_INTERPRETACION.md de esa carpeta), mas una columna de TOST que
% se agrega cuando CODIGOS/POTENCIA_EQUIVALENCIA/TOST_Core.m se corra
% sobre los datos reales. NO llenar con numeros inventados.]
Please insert text here.

\begin{table}[!t]
\caption{Tracking fidelity per subject (Comparison 3): commanded CSV
trajectory reproduced by the simulator vs. each subject's own natural
gait capture.}
\label{tab:comparacion3}
\centering
\begin{tabular}{lccccc}
\hline
Subject & RMSEnorm & $r$ & ICC(3,1) & SPM1D \% & TOST \\
\hline
[PENDIENTE] & & & & & \\
\hline
\end{tabular}
\end{table}

% \begin{figure}[!t]
% \centering
% \includegraphics[width=\columnwidth]{figuras/spm1d_comparacion3_subjectN.png}
% \caption{SPM\{t\} curve, subject N: commanded vs. own capture (see
% CODIGOS/ESTADISTICA/SPM1D_Core.m output, spm1d\_curva.png).}
% \label{fig:spm_comparacion3}
% \end{figure}

\subsection*{Representativeness of the Default Trajectory}
% [RESUELTO 13-ago-2026, esqueleto: valores de
% resultado.spm_comparacion4 (SPM1D_Core.m, diseno independiente,
% pooled trials vs. trayectoria fija) mas TOST_Core.m sobre pico/ROM.]
Please insert text here.

\begin{table}[!t]
\caption{Representativeness of the default trajectory (Comparison 4):
pooled new-subject trials vs. the simulator's fixed default trajectory.}
\label{tab:comparacion4}
\centering
\begin{tabular}{lcc}
\hline
Metric & SPM1D (\% cycle significant) & TOST \\
\hline
Platform angle (support) & [PENDIENTE] & [PENDIENTE] \\
Platform angle (swing) & [PENDIENTE] & [PENDIENTE] \\
Vertical force & [PENDIENTE] & [PENDIENTE] \\
\hline
\end{tabular}
\end{table}

\subsection*{Repeatability and Between-Subject Variability}
% [RESUELTO 13-ago-2026, esqueleto: ICC(3,1) por condicion +
% resultado.variabilidad_grupo (pico_CV_pct, ROM_CV_pct) de
% Procesar_Multisujeto_Core.m.]
Please insert text here.

\begin{table}[!t]
\caption{Between-repetition ICC(3,1) per condition and between-subject
coefficient of variation (Comparison 6).}
\label{tab:comparacion6}
\centering
\begin{tabular}{lcc}
\hline
Condition & ICC(3,1) & CV (\%) \\
\hline
[PENDIENTE] & & \\
\hline
\end{tabular}
\end{table}

\subsection*{Vertical Force: Raw, Corrected, and Literature Benchmark}
% [RESUELTO 13-ago-2026, esqueleto: fila de referencia real ya
% disponible hoy (157.3 %BW simulador vs. 98.83/104.88 %BW referencia
% del proyecto, RMSEnorm=21.1 -- ver CLAUDE.md sesion 03-ago-2026), las
% demas filas (corregida, literatura) dependen de la calibracion de
% offset, bloqueada por la integracion Raspberry Pi-ESP32.]
Please insert text here.

\begin{table}[!t]
\caption{Peak vertical force (\%BW): raw simulator output, corrected
output, and published prosthetic gait benchmarks
(see docs/literatura/literatura\_GRF\_protesica.md).}
\label{tab:fuerza}
\centering
\begin{tabular}{lcc}
\hline
Source & Peak Fz (\%BW) & RMSEnorm vs. reference \\
\hline
Simulator (raw) & 157.3 (SD 5.8) & 21.1 \\
Simulator (corrected) & [PENDIENTE] & [PENDIENTE] \\
Project reference (AMTI) & 98.83 / 104.88 & --- \\
Literature benchmark & [PENDIENTE] & --- \\
\hline
\end{tabular}
\end{table}

% ============================================================
\section{Discussion}
% ============================================================
% [PENDIENTE: interpretacion de los resultados -- depende de los datos
% de sujetos nuevos (Comparaciones 3/4/6), no se redacta con supuestos.
% INCLUIR AQUI: parrafo de traduccion clinica/TRL -- JTEHM prioriza TRL
% 5-9 y rechaza "lab validation without clinical context". Explicar la
% ruta de adopcion: quien usaria esta herramienta (protesistas,
% laboratorios de ingenieria clinica), como reduce tiempo/costo/riesgo
% de pruebas continuas con sujetos humanos durante el desarrollo de
% protesis. Ver docs/manuscrito/guia_autor_JTEHM.md seccion 1-ter para
% el hallazgo completo y precedentes publicados. La frase de
% posicionamiento TRL ya aplicada al cierre de la Introduccion (ver
% arriba) puede repetirse/ampliarse aqui si el espacio de 8 paginas lo
% permite -- medir antes de decidir (P-6, pospuesta a proposito).]
Please insert text here.

\subsection*{Limitations}
% [RESUELTO 13-ago-2026, borrador -- ver DISCUSION_Q2.md tarea de
% Discussion/Limitaciones. Redactado con lo que ya se sabe HOY, sin
% depender de datos de sujetos nuevos. Revisar y ajustar cuando lleguen
% los resultados reales -- ninguna cifra de esta subseccion se inventa,
% son las limitaciones estructurales ya declaradas en CLAUDE.md/
% DISCUSION_Q2.md, puestas en prosa por primera vez.]

Several limitations of the present design should be stated explicitly.
First, the vertical force correction combines the initial offset
calibration and the actuator tracking-fidelity check into empirically
quantified sources of error, but does not include a separately
measured per-axis inertial correction term. The moving assembly's
weight and force are not uniformly distributed across the platform, so
a single measured or CAD-derived mass value would not represent it
correctly; instead, the inertial contribution is inferred indirectly
from the same height-sweep offset calibration, which does not fully
isolate inertial from offset-related residuals. This is reported as a
limitation, not a missing term concealed as resolved.

Second, the inertial motion capture system used throughout this study
has been validated against an optoelectronic reference over an
operating range of approximately 0.8--1.6~m/s
\cite{Piche2022iSenValidity}. The simulator's own operating speed, at
approximately 30 times slower than natural human gait, falls below
this validated range; while the reported trend in that validation
favors lower speeds, applying the same accuracy figures outside the
directly tested range remains an extrapolation.

Third, cross-instrument validation --- comparing the inertial system
against an independent optoelectronic or IMU-based reference collected
in parallel on the same subjects --- was outside the scope of this
study cycle. Instrument trustworthiness instead relies on external
validation reported in the literature rather than an internal
cross-check performed by this study.

Fourth, the target sample size for the multi-subject comparisons is
supported by an a priori power analysis using the same permutation-based
statistical procedure applied in Results. That analysis approximates
between-subject variability from the trial-to-trial variability of a
single reference subject, which likely underestimates the true
variability across a population of new subjects; the analysis is
revisited once early multi-subject data become available.

Fifth, the premise motivating this study --- that the representativeness
of a single-subject-programmed trajectory across a broader population
is rarely tested with quantified, multi-subject evidence --- rests on a
limited body of directly comparable robotic gait simulator literature
rather than a systematic review confirming the gap.

\subsection*{Future Work}
% [RESUELTO 13-ago-2026, borrador -- contenido ya decidido en
% CLAUDE.md, puesto en prosa. No depende de datos.]

The present study does not modify the simulator's operating principle:
trajectories remain pre-recorded CSV files executed open-loop by the
onboard controller. Closed-loop control and algorithmic trajectory
generation, which would allow the platform to synthesize new gait
trajectories rather than replay captured ones, are reserved for a
separate, subsequent study. Likewise, an independently developed
inertial sensor and load cell built by another team member, and their
cross-validation against the instrumentation used here, were excluded
from the present scope and are planned as part of that same follow-on
work.

% ============================================================
\section{Conclusion}
% ============================================================
% [PENDIENTE: cierre breve, distinto de Discussion -- que se logro, no
% por que ni las limitaciones (eso ya va en Discussion)]

\vspace*{0.5 cm}

\section*{Acknowledgment}
% [PENDIENTE. JTEHM exige declarar fuentes de financiamiento y
% conflictos de interes aqui, y declaracion de aprobacion etica para
% investigacion con sujetos humanos.]
Please insert text here.

% Referencias ya identificadas para incluir (ver docs/literatura/):
%   - validacion_instrumentos_IMU.md: validez iSen/STT-IWS (confirmar cita exacta antes de usar)
%   - validacion_instrumentos_IMU.md: IMU en usuarios de protesis transtibial (Sensors 2023, DOI 10.3390/s23031738)
%   - literatura_GRF_protesica.md: benchmark de vGRF en marcha protesica
%   - normas_ISO_relevantes.md: ISO 5725; Wu et al. 2002 (convencion ISB)
%   - postprocesado_datos_crudos_IMU.md: filtrado/calibracion de IMU
%   - Koo & Li (2016) -- ICC(3,1)
%   - Nichols & Holmes (2002); Pataky et al. (2015) -- SPM1D no parametrico
%   - Kadaba et al. (1989) -- CMC
%   - Introduccion (nuevo 03-ago-2026, sin verificacion de texto completo
%     todavia -- ver metodos_introduccion_borrador.md):
%     "A Robotic Test Rig for Performance Assessment of Prosthetic Joints" (PMC8936071)
%     "A compact and cost-effective gait simulator to advance prosthesis
%       development with reduced reliance on human subject testing" (ScienceDirect)
%     "Conceptual design of a gait simulator for testing lower-limb active prostheses" (ResearchGate)
%     "The use of a robotic gait simulator for the development of an
%       alignment tool for lower limb prostheses" (PMC4101690)

\vspace*{1 cm}

\bibliographystyle{ieeetr}
\bibliography{references}

\end{document}
